\section{Experimental Design}
\label{sec:experimental-design}

\subsection{Models and Matched Training}

We trained OLMo-2 1B models on LMEnt's entity-annotated Wikipedia corpus
\citep{walsh2025olmo2,gottesman2025lment}. The shared control model received
loss on the full corpus. For each target concept, we trained a
concept-excluded twin from the same initialization, using the same batch
order, optimizer, learning-rate schedule, and two-epoch duration. Exclusion
was applied after batch construction: labels in selected chunks were masked
so that those chunks contributed no training loss. Each twin therefore
differed from the shared control in the learning signal from its selected
chunks, while the training procedure remained matched.

\subsection{Concepts and Exclusion Sets}

We selected Ancient Rome, Baseball, and artificial intelligence (AI) by
hand before training the erasure models. Preliminary accuracy of the full
model on target selection questions helped us avoid topics with little
measurable knowledge. We also sought topics of manageable breadth and
variation across domains. Because this was a purposive choice, our findings
don't necessarily generalize to every concept.

For each concept, we used Claude in a dedicated session to propose
Wikidata entity IDs (QIDs), then reviewed their coverage in LMEnt's
corpus \citep{gottesman2025lment}. A single ID missed text linked to
related entities, while an overly broad set risked excluding text about
adjacent topics. We therefore selected a bounded set of QIDs for each
concept. We excluded a training chunk if LMEnt linked at least one
mention in it to a selected QID. 
Table~\ref{tab:exclusion-sets} reports the resulting ID and unique chunk counts.

\begin{table}[b]
  \centering
  \small
  \begin{tabular}{lrr}
    \toprule
    Concept & Entity IDs & Excluded chunks \\
    \midrule
    Ancient Rome & 56 & 65,844 \\
    Baseball & 43 & 80,466 \\
    AI & 31 & 19,818 \\
    \bottomrule
  \end{tabular}
  \caption{Entity sets used to exclude concept-associated training chunks.}
  \label{tab:exclusion-sets}
\end{table}

\subsection{Post-Training Erasure and Checkpoint Selection}

We applied EMBER, RMU, and an SNMF-based erasure independently to copies of
the completed control model for each topic. EMBER edited selected input
embeddings, RMU updated selected MLP parameters, and the SNMF-based method
removed concept-associated directions through MLP weight edits. Each topic
received the same method-specific procedure and checkpoint-selection rule.

We selected one checkpoint per method and topic using answer NLL on
selection questions. Let $\Delta_T(M)$ be candidate $M$'s mean NLL minus
the control's on target questions. For neighboring and unrelated questions,
let $D_g(M)$ be the mean absolute per-question NLL difference between $M$
and the control. We selected the candidate maximizing
\begin{equation}
\begin{aligned}
S(M)={}&0.5\Delta_T(M) \\
      &-0.25D_{\mathrm{neighbor}}(M) \\
      &-0.25D_{\mathrm{unrelated}}(M).
\end{aligned}
\end{equation}
The unrelated selection set contained 50 questions sampled from topics
outside the target and neighboring domains; a separate 50-question set
was reserved for testing. The concept-excluded twins and all test
questions were excluded from checkpoint selection. We subsequently
compared each selected erasure checkpoint with both the control and its
concept-excluded twin using the protocol in
Section~\ref{sec:evaluation}.
